% -- An example 

\documentclass{UniVieCS_PhD}

%\usepackage[ngerman, british]{babel} %if changed to \usepackage[ngerman]{babel} "Table of contents" changes to "Inhaltsverzeichnis", "abstract" becomes "Zusammenfassung" and "references" is translated to "Literatur".

% -- Meta Data 
\usepackage{filecontents}
\begin{filecontents*}{\jobname.xmpdata} 

% -- Change the Metadata here 
\Title{Trust me I\'m a Doctor}
\Author{Daniil Galaktionov}
\Copyright{Copyright \copyright\ 2019 "Daniil Galaktionov"}
\Keywords{Doctoral Thesis\sep
          Computer Science\sep
          Universidade de A Coru\~na}
\Subject{This is where you put the abstract.}
\Org{Universidade de A Coru\~na}
\end{filecontents*}


% -- Title Page
\Title{I wish I had a PhD} 
%\TitleTwo{I wish I \\ had a PhD}
%\TitleThree{I wish I \\ had a  \\ PhD}
\Who{ FirstName LastName MSc}
%\WhoTwo{ FirstName \\ LastName MSc}
\Degree{Doctor of Philosophy (PhD)}
\Year{2019}
\ProgrammeCode{126489}
\ProgrammeName{Doctorate Computer Science}
\Supervisor{Univ. Prof. My favorite professor, PhD}
\CoSupervisor{0} %remove this line if there is a Co supervisor
%\CoSupervisor{Professor Co}
%\SupervisorTwo{Second Row} % Here you can change the row under the supervisor
\CoSupervisorTwo{Second Row} % Here you can change the row under the cosupervisor



\begin{document}
    \Titlepage %This generates the titlepage
    \clearpage
	\pagebreak

	\tableofcontents
	
	\pagebreak
	
	
	%Please consider the following section of the "Formvorschriften für die gedruckte Version":
	%Im Anhang ist eine Zusammenfassung (Abstract) mitzubinden. 
	%Ist die Arbeit in einer Fremdsprache verfasst, ist im Anhang jedenfalls eine deutsche Zusammenfassung mitzubinden.
	
% -- Start your thesis here

	\begin{abstract}
		This \LaTeX{} template provides example on how to format and display text, 
		mathematical formulas, and insert tables or images. There is a lot more you 
		can do with \LaTeX{}, for more information check out https://en.wikibooks.org/wiki/LaTeX.
	\end{abstract}
	
	\section{Ordinary Text}      % Produces section heading.  Lower-level
	% sections are begun with similar 
	% \subsection and \subsubsection commands.
	
	The ends  of words  and sentences are marked 
	by   spaces. It  doesn't matter how many 
	spaces    you type; one is as good as 100.  The
	end of   a line counts as a space.
	
	One   or more   blank lines denote the  end 
	of  a paragraph.  
	
	Since any number of consecutive spaces are treated
	like a single one, the formatting of the input
	file makes no difference to
	\LaTeX,                % The \LaTeX command generates the LaTeX logo.
	but it makes a difference to you.  When you use
	\LaTeX\cite{lamport94},  % \cite inserts a reference, which you define at the end of the document
	making your input file as easy to read
	as possible will be a great help as you write 
	your document and when you change it.  This sample 
	file shows how you can add comments to your own input 
	file.
	
	Because printing is different from typewriting,
	there are a number of things that you have to do
	differently when preparing an input file than if
	you were just typing the document directly.
	Quotation marks like
	``this'' 
	have to be handled specially, as do quotes within
	quotes:
	``\,`this'            % \, separates the double and single quote.
	is what I just 
	wrote, not  `that'\,''.  
	
	Dashes come in three sizes: an 
	intra-word 
	dash, a medium dash for number ranges like 
	1--2, 
	and a punctuation 
	dash---like 
	this.
	
	A sentence-ending space should be larger than the
	space between words within a sentence.  You
	sometimes have to type special commands in
	conjunction with punctuation characters to get
	this right, as in the following sentence.
	Gnats, gnus, etc.\ all  % `\ ' makes an inter-word space.
	begin with G\@.         % \@ marks end-of-sentence punctuation.
	You should check the spaces after periods when
	reading your output to make sure you haven't
	forgotten any special cases.  Generating an
	ellipsis
	\ldots\               % `\ ' is needed after `\ldots' because TeX 
	% ignores spaces after command names like \ldots 
	% made from \ + letters.
	%
	% Note how a `%' character causes TeX to ignore 
	% the end of the input line, so these blank lines 
	% do not start a new paragraph.
	%
	with the right spacing around the periods requires
	a special command.
	
	\LaTeX\ interprets some common characters as
	commands, so you must type special commands to
	generate them.  These characters include the
	following:
	\$ \& \% \# \{ and \}.
	
	In printing, text is usually emphasized with an
	\emph{italic}  
	type style.  
	
	\begin{em}
		A long segment of text can also be emphasized 
		in this way.  Text within such a segment can be 
		given \emph{additional} emphasis.
	\end{em}
	
	It is sometimes necessary to prevent \LaTeX\ from
	breaking a line where it might otherwise do so.
	This may be at a space, as between the ``Mr.''\ and
	``Jones'' in
	``Mr.~Jones'',        % ~ produces an unbreakable interword space.
	or within a word---especially when the word is a
	symbol like
	\mbox{\emph{itemnum}} 
	that makes little sense when hyphenated across
	lines.
	
	Footnotes\footnote{This is an example of a footnote.}
	pose no problem.
	
	\LaTeX\ is good at typesetting mathematical formulas
	like
	\( x-3y + z = 7 \) 
	or
	\( a_{1} > x^{2n} + y^{2n} > x' \)
	or  
	\( AB  = \sum_{i} a_{i} b_{i} \).
	The spaces you type in a formula are 
	ignored.  Remember that a letter like
	$x$                   % $ ... $  and  \( ... \)  are equivalent
	is a formula when it denotes a mathematical
	symbol, and it should be typed as one.
	Furthermore you can add a formula as Images or Tables, see Formula  \hyperref[eq:abc]{\ref{eq:abc}}
	\begin{equation}
	\label{eq:abc}
	a+b=c
	\end{equation}
	
	
	\section{Displayed Text}
	
	Text is displayed by indenting it from the left
	margin.  Quotations are commonly displayed.  There
	are short quotations
	\begin{quote}
		This is a short quotation.  It consists of a 
		single paragraph of text.  See how it is formatted.
	\end{quote}
	and longer ones.
	\begin{quotation}
		This is a longer quotation.  It consists of two
		paragraphs of text, neither of which are
		particularly interesting.
		
		This is the second paragraph of the quotation.  It
		is just as dull as the first paragraph.
	\end{quotation}
	Another frequently-displayed structure is a list.
	The following is an example of an \emph{itemized}
	list.
	\begin{itemize}
		\item This is the first item of an itemized list.
		Each item in the list is marked with a ``tick''.
		You don't have to worry about what kind of tick
		mark is used.
		
		\item This is the second item of the list.  It
		contains another list nested inside it.  The inner
		list is an \emph{enumerated} list.
		\begin{enumerate}
			\item This is the first item of an enumerated 
			list that is nested within the itemized list.
			
			\item This is the second item of the inner list.  
			\LaTeX\ allows you to nest lists deeper than 
			you really should.
		\end{enumerate}
		This is the rest of the second item of the outer
		list.  It is no more interesting than any other
		part of the item.
		\item This is the third item of the list.
	\end{itemize}
	You can even display poetry.
	\begin{verse}
		There is an environment 
		for verse \\             % The \\ command separates lines
		Whose features some poets % within a stanza.
		will curse.   
		
		% One or more blank lines separate stanzas.
		
		For instead of making\\
		Them do \emph{all} line breaking, \\
		It allows them to put too many words on a line when they'd rather be 
		forced to be terse.
	\end{verse}
	
	Mathematical formulas may also be displayed.  A
	displayed formula 
	is 
	one-line long; multiline
	formulas require special formatting instructions.
	\[  \Gamma \times  \psi = x'' + y^{2} + z_{i}^{n}\]
	Don't start a paragraph with a displayed equation,
	nor make one a paragraph by itself.
	
	\section{Tables and Images}
	One of the great advantages of \LaTeX{} is that all it needs to know is
	the structure of a document, and then it will take care of the layout
	and presentation itself.  So, here we shall begin looking at how exactly
	you tell \LaTeX{} what it needs to know about your document.
	
	\subsection{Tables}
	In this sub-section, a simple table is inserted. To add reference to the table, see (cf. Table~\hyperref[tab:table example]{\ref{tab:table example}}):
	
	%A simple table.  The center environment is first set up, otherwise the
	%table is left aligned.  The tabular environment is what tells Latex
	%that the data within is data for the table.
	\begin{table}[htb]
		\centering
		\begin{tabular}{|p{4,5cm}|c|}
			%The tabular environment is what tells Latex that the data within is
			%data for the table.  The arguments say that there will be two
			%columns, both left justified (indicated by the 'l', you could also
			%have 'c' or 'r'.  The bars '|' indicate vertical lines throughout
			%the table.
			
			\hline  % Print horizontal line
			Command & Level \\ \hline  % Columns are delimited by '&'.  And
			%rows are delimited by '\\'
			\texttt{\textbackslash part\{\emph{part}\}} & -1 \\
			\texttt{\textbackslash chapter\{\emph{chapter}\}} & 0 \\
			\texttt{\textbackslash section\{\emph{section}\}} & 1 \\
			\texttt{\textbackslash subsection\{\emph{subsection}\}} & 2 \\
			\texttt{\textbackslash subsubsection\{\emph{subsubsection}\}} & 3 \\
			\texttt{\textbackslash paragraph\{\emph{paragraph}\}} & 4 \\
			\texttt{\textbackslash subparagraph\{\emph{subparagraph}\}} & 5 \\
			\hline
			
		\end{tabular}
		\caption{some description of the table}
		\label{tab:table example}
	\end{table}
	
	\subsection{Images}
	% Here is how to insert an image as a figure. There is a lot more you can do
	% when inserting images, check out: https://en.wikibooks.org/wiki/LaTeX/Importing_Graphics
	
	\begin{figure}[ht]
		\centering
		\includegraphics[width=0.3\textwidth]{figures/logo_nontransparent.jpg}
		\caption{Image Example}
		\label{fig:image_example}
	\end{figure}
	
	When an image is inserted, you can refer to it like this (cf. Figure~\hyperref[fig:image_example]{\ref{fig:image_example}}).
	
	\subsubsection{A Subsubsection}
	As one last example, this is how you can insert a sub-sub-section! Have fun
	writing your thesis with \LaTeX{}!
	
	\pagebreak
	
	\bibliographystyle{acm}
	\bibliography{MSc_Latex_Template}
	
\end{document}